
\documentclass{article}
\usepackage{colortbl}
\usepackage{makecell}
\usepackage{multirow}
\usepackage{supertabular}

\begin{document}

\newcounter{utterance}

\centering \large Interaction Transcript for game `eqbench', experiment `eq\_bench', episode 55 with reflexion{-}llama{-}3.1{-}8b{-}instruct.
\vspace{24pt}

{ \footnotesize  \setcounter{utterance}{1}
\setlength{\tabcolsep}{0pt}
\begin{supertabular}{c@{$\;$}|p{.15\linewidth}@{}p{.15\linewidth}p{.15\linewidth}p{.15\linewidth}p{.15\linewidth}p{.15\linewidth}}
    \# & \multicolumn{2}{c}{Player} && \multicolumn{2}{c}{Game Master} \\
    \hline

    \theutterance \stepcounter{utterance}  
    & & & \multicolumn{4}{p{0.6\linewidth}}{
        \cellcolor[rgb]{0.9,0.9,0.9}{
            \makecell[{{p{\linewidth}}}]{
                \texttt{\tiny{[P1$\langle$GM]}}
                \texttt{You have attempted this same task before and did not succeed. These are the plans you wrote after those attempts.} \\
\texttt{Trial \#1: New plan: } \\
\\ 
\texttt{Upon re{-}examining the dialogue, I realize that I focused too much on the surface{-}level exchange between Reggie and Marcus, without considering the underlying dynamics of their relationship. I should have taken into account the fact that Marcus is Reggie's sibling, and their banter often involves a mix of sarcasm and affection.} \\
\\ 
\texttt{To improve my analysis, I will pay closer attention to the tone and language used by each character, particularly in their interactions with each other. I will also consider the context of the conversation, including Reggie's idea for a birthday party and Marcus's initial skepticism.} \\
\\ 
\texttt{In this case, I should have scored Apologetic higher, as Marcus's response to Reggie's jab about forgetting to show up suggests that he is not entirely comfortable with the direction of the conversation. I will also consider the possibility that Marcus may be feeling a mix of emotions, including Provoked and Apologetic.} \\
\\ 
\texttt{Defeated: 0} \\
\texttt{Provoked: 8} \\
\texttt{Apologetic: 6} \\
\texttt{Indifferent: 0} \\
\\ 
\texttt{Trial \#2: New plan: } \\
\\ 
\texttt{Upon re{-}examining the dialogue, I realize that I still focused too much on the surface{-}level exchange between Reggie and Marcus, without considering the underlying dynamics of their relationship. I should have taken into account the fact that Marcus is Reggie's sibling, and their banter often involves a mix of sarcasm and affection.} \\
\\ 
\texttt{To improve my analysis, I will pay closer attention to the tone and language used by each character, particularly in their interactions with each other. I will also consider the context of the conversation, including Reggie's idea for a birthday party and Marcus's initial skepticism.} \\
\\ 
\texttt{However, I also realize that I should have considered the emotional impact of Reggie's jab about Marcus forgetting to show up. This comment is likely to have stung Marcus, and he may be feeling a mix of emotions, including Provoked, Apologetic, and possibly even a hint of Defeated.} \\
\\ 
\texttt{To improve my analysis, I will also consider the possibility that Reggie's comment may have been a deliberate attempt to provoke Marcus, and that Reggie may be feeling a sense of satisfaction or even glee at having gotten a rise out of his sibling.} \\
\\ 
\texttt{Defeated: 2} \\
\texttt{Provoked: 9} \\
\texttt{Apologetic: 7} \\
\texttt{Indifferent: 0} \\
\\ 
\\ 
\\ 
\texttt{Your task is to predict the likely emotional responses of a character in this dialogue:} \\
\\ 
\texttt{Reggie: Marcus, I've been thinking about Mom's birthday party. I've got some ideas.} \\
\texttt{Marcus: Oh, I bet you do. Another corporate gala, Reggie? } \\
\texttt{Reggie: Funny. I was actually thinking of something more... intimate. } \\
\texttt{Marcus: Intimate? You? Now, that's a surprise.} \\
\texttt{Reggie: Yes, well, I can be full of surprises. } \\
\texttt{Marcus: I'll believe it when I see it. } \\
\texttt{Reggie: What's that supposed to mean? } \\
\texttt{Marcus: Nothing. Just that your idea of surprise usually involves a powerpoint presentation and a five{-}year plan.} \\
\texttt{Reggie: And your idea of a surprise is forgetting to show up. } \\
\texttt{Marcus: Low blow, sis. } \\
\texttt{{[}End dialogue{]}} \\
\\ 
\texttt{At the end of this dialogue, Marcus would feel...} \\
\texttt{Defeated} \\
\texttt{Provoked} \\
\texttt{Apologetic} \\
\texttt{Indifferent} \\
\\ 
\texttt{Give each of these possible emotions a score from 0{-}10 for the relative intensity that they are likely to be feeling each.} \\
\\ 
\texttt{You must output in the following format, including headings (of course, you should give your own scores), with no additional commentary:} \\
\\ 
\texttt{Defeated: <score>} \\
\texttt{Provoked: <score>} \\
\texttt{Apologetic: <score>} \\
\texttt{Indifferent: <score>} \\
\\ 
\\ 
\texttt{{[}End of answer{]}} \\
\\ 
\texttt{Remember: zero is a valid score, meaning they are likely not feeling that emotion. You must score at least one emotion > 0.} \\
\\ 
\texttt{Your answer:} \\
            }
        }
    }
     \\ \\

    \theutterance \stepcounter{utterance}  
    & \multicolumn{4}{p{0.6\linewidth}}{
        \cellcolor[rgb]{0.561,0.737,0.561}{
            \makecell[{{p{\linewidth}}}]{
                \texttt{\tiny{[P1$\rangle$GM]}}
                \texttt{Defeated: 4} \\
\texttt{Provoked: 9} \\
\texttt{Apologetic: 2} \\
\texttt{Indifferent: 0} \\
            }
        }
    }
    & & \\ \\

    \theutterance \stepcounter{utterance}  
    & & & \multicolumn{2}{p{0.3\linewidth}}{
        \cellcolor[rgb]{0.9,0.9,0.9}{
            \makecell[{{p{\linewidth}}}]{
                \texttt{\tiny{[GM$|$GM]}}
                \texttt{Defeated: 4} \\
\texttt{Provoked: 9} \\
\texttt{Apologetic: 2} \\
\texttt{Indifferent: 0} \\
            }
        }
    }
    & & \\ \\

    \theutterance \stepcounter{utterance}  
    & & & \multicolumn{2}{p{0.3\linewidth}}{
        \cellcolor[rgb]{0.9,0.9,0.9}{
            \makecell[{{p{\linewidth}}}]{
                \texttt{\tiny{[GM$|$GM]}}
                \texttt{Defeated: 0} \\
\texttt{Provoked: 7} \\
\texttt{Apologetic: 0} \\
\texttt{Indifferent: 0} \\
            }
        }
    }
    & & \\ \\

    \theutterance \stepcounter{utterance}  
    & & & \multicolumn{2}{p{0.3\linewidth}}{
        \cellcolor[rgb]{0.9,0.9,0.9}{
            \makecell[{{p{\linewidth}}}]{
                \texttt{\tiny{[GM$|$GM]}}
                \texttt{game\_result = LOSE} \\
            }
        }
    }
    & & \\ \\

    \theutterance \stepcounter{utterance}  
    & & \multicolumn{4}{p{0.6\linewidth}}{
        \cellcolor[rgb]{0.95,0.95,0.95}{
            \makecell[{{p{\linewidth}}}]{
                \texttt{\tiny{[P1$\leftrightarrow$P]}}
                \texttt{Post{-}episode Reflexion (not sent to GM):} \\
\texttt{New plan: } \\
\\ 
\texttt{Upon re{-}examining the dialogue, I realize that I still focused too much on the surface{-}level exchange between Reggie and Marcus, without considering the underlying dynamics of their relationship. I should have taken into account the fact that Marcus is Reggie's sibling, and their banter often involves a mix of sarcasm and affection.} \\
\\ 
\texttt{To improve my analysis, I will pay closer attention to the tone and language used by each character, particularly in their interactions with each other. I will also consider the context of the conversation, including Reggie's idea for a birthday party and Marcus's initial skepticism.} \\
\\ 
\texttt{However, I also realize that I should have considered the emotional impact of Reggie's jab about Marcus forgetting to show up. This comment is likely to have stung Marcus, and he may be feeling a mix of emotions, including Provoked, Apologetic, and possibly even a hint of Defeated.} \\
\\ 
\texttt{To improve my analysis, I will also consider the possibility that Reggie's comment may have been a deliberate attempt to provoke Marcus, and that Reggie may be feeling a sense of satisfaction or even glee at having gotten a rise out of his sibling.} \\
\\ 
\texttt{Additionally, I will consider the fact that Reggie's comment about Marcus's idea of a surprise being a PowerPoint presentation and a five{-}year plan is a clear jab at Marcus's corporate job, and that Reggie may be feeling a sense of superiority or disdain for Marcus's career choices.} \\
\\ 
\texttt{To improve my analysis, I will also consider the possibility that Marcus may be feeling a sense of defensiveness or hurt pride in response to Reggie's comment, and that he may be struggling to navigate his feelings about his job and his relationship with Reggie.} \\
\\ 
\texttt{Defeated: 3} \\
\texttt{Provoked: 9} \\
\texttt{Apologetic: 6} \\
\texttt{Indifferent: 0} \\
            }
        }
    }
    & \\ \\

\end{supertabular}
}

\end{document}
